\chapter{Tics}

\section*{Definition}
Tics are sudden, rapid, recurrent, non-rhythmic motor movements or vocalizations that are involuntary but often preceded by a premonitory urge, are temporarily suppressible, and may be briefly mimicked voluntarily. Tic disorders are classified by duration and type into transient tic disorder, chronic (persistent) motor or vocal tic disorder, and Tourette syndrome (multiple motor and at least one vocal tic for more than one year).

\section*{Pathophysiology}
Tics involve dysfunction of cortico-striato-thalamo-cortical circuits, particularly the sensorimotor and prefrontal loops through the caudate and putamen. Reduced inhibition within the striatum, especially decreased GABAergic interneuron activity, leads to excessive disinhibition of thalamocortical motor pathways. Dopamine dysregulation plays a central role, with evidence of both presynaptic dopamine release and postsynaptic receptor hypersensitivity. Genetic factors are prominent, with linkage to SLITRK1, HDC, and multiple common variants with small effect sizes.

\section*{Etiology}
\begin{itemize}
  \item Primary tic disorders: Tourette syndrome, chronic motor or vocal tic disorder, transient tic disorder (most common, up to 20\% of school-age children)
  \item Genetic factors: Family history common; monogenic forms rare (SLITRK1, HDC, CNTNAP2, NRXN1); polygenic heritability estimates 30--80\%
  \item Associated conditions: Obsessive-compulsive disorder (OCD), attention-deficit/hyperactivity disorder (ADHD), anxiety disorders, autism spectrum disorder, self-injurious behaviors
  \item Secondary (tic-like) causes: Drug-induced (cocaine, amphetamines, methylphenidate, antipsychotics, levodopa, anticonvulsants), head trauma, encephalitis (herpes simplex, anti-NMDA receptor encephalitis), stroke, carbon monoxide poisoning, neuroacanthocytosis, Huntington disease, Sydenham chorea (PANDAS pediatric autoimmune neuropsychiatric disorders)
  \item Functional tic-like movements: Abrupt onset in adolescence/young adulthood, often with complex whole-body movements and vocalizations, not preceded by premonitory urge, and associated with social media exposure
\end{itemize}

\section*{Indications for Evaluation}
Seek care if:
\begin{itemize}
  \item Tics causing functional impairment, social distress, or physical injury
  \item Onset after age 15 (atypical for primary tic disorders, consider secondary or functional causes)
  \item Rapidly progressive or explosive symptom escalation
  \item Associated symptoms suggesting underlying neurologic disorder (seizures, cognitive decline, chorea, weakness)
  \item Severe or persistent tics despite behavioral therapy
\end{itemize}

\section*{Related Symptoms}
\begin{itemize}
  \item Premonitory urge
  \item Motor tics (simple: eye blinking, head jerking; complex: hopping, touching, echopraxia, copropraxia)
  \item Vocal tics (simple: throat clearing, grunting; complex: echolalia, palilalia, coprolalia)
  \item Obsessions and compulsions
  \item Inattention and hyperactivity
  \item Anxiety
  \item Self-injurious behaviors
  \item Sleep disturbances
\end{itemize}
\textbf{Note:} Coprolalia (involuntary swearing) occurs in only 10--15\% of Tourette syndrome patients despite popular perception; the premonitory urge preceding tics distinguishes them from other hyperkinetic movement disorders.

\section*{Self-Care / Home Management}
\begin{itemize}
  \item Educate the patient, family, and teachers about the involuntary nature of tics to reduce stigma and punishment.
  \item Avoid scolding or drawing attention to tics as this may increase anxiety and worsen symptoms.
  \item Encourage activities that improve concentration and reduce stress such as exercise, music, and hobbies.
  \item Ensure adequate sleep, as fatigue significantly exacerbates tics.
\end{itemize}

\section*{Diagnostic Approach}
\begin{itemize}
  \item Clinical diagnosis based on DSM-5 criteria: onset before age 18, multiple motor and/or vocal tics, duration of symptoms > 1 year for Tourette syndrome.
  \item History from patient and caregivers including tic types, age of onset, premonitory urges, suppressibility, exacerbating factors, and family history.
  \item Screening for comorbid conditions (ADHD, OCD, anxiety) using validated tools (Yale Global Tic Severity Scale, Y-BOCS).
  \item Secondary causes are considered if atypical features present (onset after age 15, explosive progression, abnormal neurologic exam); brain MRI and laboratory testing reserved for these cases.
\end{itemize}

\section*{Treatment / Management Overview}
\begin{itemize}
  \item Behavioral therapy: Comprehensive Behavioral Intervention for Tics (CBIT), including habit reversal training and exposure with response prevention, is first-line treatment.
  \item Pharmacotherapy for moderate to severe tics: Alpha-2 agonists (clonidine, guanfacine) are first-line; antipsychotics (risperidone, aripiprazole, haloperidol, pimozide) are reserved for refractory cases.
  \item Treat comorbid disorders: Stimulants or atomoxetine for ADHD, SSRIs for OCD, benzodiazepines for anxiety.
  \item Deep brain stimulation (thalamic or pallidal) is reserved for severely disabling, medication-refractory Tourette syndrome in adults.
  \item Education and reassurance for mild tics causing no functional impairment.
\end{itemize}

\clearpage
